\section{Stabilizers and parabolic intersections}
\label{sec:stabilizers}

We now combine the two containment arguments. Throughout this section,
$X\subseteq S$ and its relative family are fixed. The element $x\in A$
is arbitrary; in particular, it need not belong to the subgroup that
is currently acting.

\begin{lemma}[Conjugating a stabilizer]
\label{stab:conjugation}
Let $D\subseteq Y\subseteq S$ and $b\in A_Y$. If
$H(Y,x)\subseteq b^{-1}A_Db$, then
\[
 bH(Y,x)b^{-1}=H(D,bx).
\]
\end{lemma}
\begin{proof}
For $c\in A_D$, invertibility of $B_b$ gives
\[
 E_L(cbx)\simeq E_L(bx)\quad\text{for all }L
 \quad\Longleftrightarrow\quad
 E_L(b^{-1}cbx)\simeq E_L(x)\quad\text{for all }L.
\]
Since $b^{-1}cb\in A_Y$, the right side says that
$b^{-1}cb\in H(Y,x)$. Thus
$H(D,bx)=A_D\cap bH(Y,x)b^{-1}$. The assumed containment removes the
intersection.
\end{proof}

\begin{theorem}[Relative stabilizers are parabolic]
\label{stab:parabolic}
For every $Y\subseteq S$ and $x\in A$, there exist $R\subseteq Y$ and
$c\in A_Y$ such that
\[
 H(Y,x)=c^{-1}A_Rc.
\]
\end{theorem}
\begin{proof}
We induct on $|Y|$, allowing every $x\in A$ at every rank. The assertion
is immediate for $Y=\varnothing$. By Lemma~\ref{stab:conjugation}, it
suffices to place $H(Y,x)$ inside $b^{-1}A_Db$ for some $b\in A_Y$ and
proper $D\subsetneq Y$: the conjugated stabilizer is then $H(D,bx)$,
to which induction applies.

Suppose first that $Y$ is spherical. If all its tests on the family at
$x$ are acyclic, every generator in $Y$ fixes every member, so
$H(Y,x)=A_Y$. Otherwise Proposition~\ref{sph:containment} gives the
required proper-type containment.

Now suppose that $Y$ is nonspherical. Choose a minimum vertex
$P=(T,A_Tz)$ of the relative harmonic filtration of $A_Yx$. Let $K$ be
its invisible set and put $T_+=T\setminus K$, as in
Section~\ref{sec:heights}. Proposition~\ref{ht:containment} gives
\[
 x^{-1}H(Y,x)x\subseteq z^{-1}A_Dz,
 \qquad D=T_+\cup K\subseteq Y.
\]
Since $z\in A_Yx$, the element $b=zx^{-1}$ belongs to $A_Y$, and this
is exactly $H(Y,x)\subseteq b^{-1}A_Db$. If $D\ne Y$, the induction
step is complete.

It remains to treat $D=Y$. The sets $T_+$ and $K$ are disjoint and
nonadjacent, and $T_+$ is spherical because it is a subset of $T$.
We claim that the entire $K$-factor fixes the family at $x$.

Let $j\notin Y$ be an exterior neighbor of $K$. At every recorded level,
the definition of $K$ makes its harmonic coordinate zero. Outside $Y$
these coordinates are actual generator counts, fixed throughout $A_Yx$.
The cutoff was chosen at or below \emph{every} occupied exterior level,
so there are no generators of color $j$ at any level, for any member
of the family at any $w\in A_Yx$. There is also no
edge from $K$ to $T_+=Y\setminus K$. If a member $E_L(w)$ had a generator
in $K$, its connected color support, which contains $f\notin Y$, would
have to cross an exterior neighbor of $K$. We have just excluded every
such generator. This applies in particular at $w=x$: the family there
has no generators in $K$ or any of its neighbors, even at levels below
the cutoff.

Every test at a color of $K$ is therefore zero termwise. The evaluation
cone shows that its positive generator fixes each member, and its inverse
does so as well. The Artin presentation and standard embedding give
\[
 A_Y=A_K\times A_{T_+},
 \qquad H(Y,x)=A_K\times H(T_+,x).
\]
Since $Y$ is nonspherical, the spherical subset $T_+$ is proper.
Induction supplies $R\subseteq T_+$ and $c\in A_{T_+}$ with
$H(T_+,x)=c^{-1}A_Rc$. Commutation with $A_K$ now gives
\[
 H(Y,x)=c^{-1}(A_K\times A_R)c=c^{-1}A_{K\cup R}c.
\]
This also covers $T_+=\varnothing$.
\end{proof}

\begin{proposition}[The initial stabilizer]
\label{stab:initial}
The relative family associated with $X$ satisfies $H(S,1)=A_X$.
\end{proposition}
\begin{proof}
Every color of $X$ is nonadjacent to $f$. Its evaluation on every
$P_{f,L}$ is zero, so its twist fixes that projective. This proves
$A_X\subseteq H(S,1)$.

Consider the filtration for $Y=S$ and $x=1$. The empty-type vertex
$P_0=(\varnothing,\{1\})$ has only one nonzero layer vector, namely
the unit vector at $f$ in level zero. This vertex has minimum height.
Indeed a competitor with any nonzero higher block is later. At level
zero, every competitor retains the exterior coordinate $p_f^0=1$, and
all coordinates are nonnegative. Its coordinate sum is therefore at
least one. Equality forces its entire vector to be the unit vector at
$f$, and then every closed-type sign coordinate is zero, since the
type lies in $S$. Every lower block of $P_0$ is zero and is least.

At this minimum $T_+=\varnothing$. A color of $S$ is invisible precisely
when it is nonadjacent to $f$, so $K=X$. Applying
Proposition~\ref{ht:containment} with $z=x=1$ gives
$H(S,1)\subseteq A_X$.
\end{proof}

\begin{proof}[Proof of Theorem~\ref{thm:main}]
Simultaneous conjugation reduces the assertion to the intersection
$A_X\cap x^{-1}A_Yx$, with $x\in A$. For $a\in A_Y$, invertibility
of $B_x$ and Proposition~\ref{stab:initial} give
\begin{align*}
 x^{-1}ax\in A_X
 &\ \Longleftrightarrow\ 
 E_L(x^{-1}ax)\simeq E_L(1)\quad\text{for every }L\\
 &\ \Longleftrightarrow\ 
 E_L(ax)\simeq E_L(x)\quad\text{for every }L\\
 &\ \Longleftrightarrow\ a\in H(Y,x).
\end{align*}
Consequently
\[
 A_X\cap x^{-1}A_Yx=x^{-1}H(Y,x)x.
\]
By Theorem~\ref{stab:parabolic}, write $H(Y,x)=c^{-1}A_Zc$ with
$Z\subseteq Y$ and $c\in A_Y$. The displayed intersection is then
$(cx)^{-1}A_Z(cx)$, hence is parabolic. Conjugating back proves the
assertion for the originally specified $g$ and $h$.
\end{proof}


\subsection{Parabolic closure and dynamics}
\label{sec:consequences}

The inside-parabolic description obtained above makes the rank decrease
in an intersection explicit. We use it to prove the introductory
consequences, then record the fixed-point and continuity implications.

\begin{proof}[Proof of Corollary~\ref{cor:arbitrary-intersections}]
If $Q=gA_Yg^{-1}$ and $P=hA_Xh^{-1}$, put $x=g^{-1}h$. The proof of
Theorem~\ref{thm:main} and Theorem~\ref{stab:parabolic} give
\[
 Q\cap P=g c^{-1}A_Rc g^{-1}
 \quad\text{for some }R\subseteq Y\text{ and }c\in A_Y.
\]
If this intersection is strictly smaller than $Q$, then $R\ne Y$, since
conjugation by $c\in A_Y$ preserves $A_Y$. Thus a strict intersection step
decreases the size of the displayed defining set.

Start with $Q=A=A_S$ and no selected members of $\mathcal F$. Whenever $Q$
is not contained in some $P\in\mathcal F$, replace $Q$ by $Q\cap P$ and
select that member. The preceding strict decrease permits at most $|S|$
steps. The resulting $Q$ is contained in every member of $\mathcal F$ and
is the intersection of the selected subfamily, proving the displayed
equality and parabolicity. Finally, the family of parabolics containing
$E$ is nonempty because it contains $A$; its intersection is the unique
smallest such subgroup.
\end{proof}

\begin{proof}[Proof of Corollary~\ref{cor:nonclique-dynamics}]
In the presentation graph in which every finite label is an edge, the
infinite-label pair gives the visual splitting over
$A_{S\setminus\{s,t\}}$ of
\cite[Definition~2.4]{CharneyMartinMorrisWright2026}.
Theorem~\ref{thm:main} makes intersections of two conjugates of this edge
group parabolic in $A$. Hence
\cite[Theorem~3.2]{CharneyMartinMorrisWright2026} gives acylindrical
hyperbolicity and weak malnormality of the edge group, and
\cite[Proposition~4.2]{CharneyMartinMorrisWright2026} gives weak malnormality
of every proper parabolic.

The finite intersection supplied by weak malnormality is parabolic by
Theorem~\ref{thm:main}. A nontrivial parabolic contains a conjugate of
some $\sigma_s$, which has infinite order: the exponent-sum homomorphism
$A\to\mathbb Z$ sends every standard generator to $1$. Thus a finite
parabolic is trivial, giving the stated witness.
\end{proof}

\subsection{Fixed-point subgroups and locally compact groups}
\label{sec:fixed-point-consequences}

A parabolic subgroup is \emph{complete} if it is conjugate to $A_T$
for a set $T$ with no infinite label; it need not be spherical.
A group has property $\mathrm{FA}'$ if, in every simplicial action on a
tree without inversions, each element fixes a vertex. A locally compact
Hausdorff group $L$ is \emph{almost connected} if $L/L^\circ$ is compact.
A discrete group is \emph{lcH-slender} if every abstract homomorphism
into it from a locally compact Hausdorff group is continuous.

\begin{corollary}[Reduction to complete parabolics]
\label{cor:complete-parabolics}
Let $A=A_M$ with $S$ finite.
\begin{enumerate}
\item Every subgroup of $A$ with property $\mathrm{FA}'$ lies in a
complete parabolic subgroup.
\item Every abstract homomorphism $L\to A$ from an almost-connected
locally compact Hausdorff group has image in a complete parabolic.
\item If every complete standard parabolic of $A$ is lcH-slender,
then $A$ is lcH-slender.
\end{enumerate}
\end{corollary}
\begin{proof}
Theorem~\ref{thm:main} places $A$ in the class of groups whose parabolics
are closed under pairwise intersections. Apply
\cite[Proposition~1.3 and Theorem~1.7]{MoellerParisVarghese2024}.
\end{proof}


\subsection{Faithfulness and effective membership}
\label{sec:algorithms}

The unit-label projective already detects membership in $A_X$. Indeed,
label compatibility in Proposition~\ref{cat:family} and the initial
stabilizer calculation give
\begin{equation}
 B_wP_{f,\mathbf1}\simeq P_{f,\mathbf1}
 \quad\Longleftrightarrow\quad w\in A_X.
 \label{stab:single-detector}
\end{equation}
For the forward implication, tensoring labels by each simple $L$ gives
$B_wP_{f,L}\simeq P_{f,L}$, so Proposition~\ref{stab:initial} applies.
The reverse implication follows from the same proposition.
In particular, for $X=\varnothing$ the constructed action is faithful:
an element acting by the identity functor fixes $P_{f,\mathbf1}$ and
is therefore the identity of $A$. This is the action for our framed,
per-edge-labelled algebra; the faithfulness question for the original
unframed construction of \cite[Conjecture~7.10]{HengLicata2024} concerns
a different category.

To make \eqref{stab:single-detector} an algorithm, we must compute in the
chosen algebra exactly. The related twist-group algorithm of
\cite[Proposition~7.11]{HengLicata2024} also uses Gaussian elimination.
We give the coefficient-field construction here, so finite-dimensionality
over $\mathbb C$ is not being used as an effectiveness assumption.

\begin{proof}[Proof of Corollary~\ref{cor:membership}]
Fix the supplied $M$ and $X$, and use their relative construction.
Let $N$ be the least common multiple of the integers $2m$ for finite
edge labels $m\ge4$, and take $K=\mathbb Q(\exp(2\pi i/N))$; take
$K=\mathbb Q$ if there is no such label. Then $K$ contains every
parameter $q_m=\exp(\pi i/m)$. This field and its exact arithmetic are
computable from the finite matrix $M$.

We realize each Temperley--Lieb--Jones factor in its planar model.
Loop-free pairing diagrams give finite bases for each required underlying
planar morphism space; stacking diagrams evaluates every closed loop as $q_m+q_m^{-1}$.
The Jones--Wenzl recursion computes the simple idempotents
$p_0,\ldots,p_{m-2}$ over $K$, with nonzero denominators in this range
\cite[Definition~2.1.1, Remark~2.1.3, and Theorem~2.3.2]{Chen2014}.
Negligible maps are kernels of finite trace-pairing matrices, so their
quotient spaces and basis reductions are obtained by exact linear
algebra. For idempotents $p,q$, their corner morphism space is the image
of $f\mapsto qfp$ in that quotient
\cite[Definition~2.4.1, Lemma~2.4.3, and Definition~5.2.1]{Chen2014}.
Composition, tensoring and pivotal evaluation are consequently
computable in these finite diagram spaces. These finite kernels,
quotients and corner images commute with scalar extension to
$\mathbb C$, recovering the model used in Section~\ref{sec:category}.

The known simple list also makes decomposition effective. For any
object $U$ needed in a product, compute bases of
$\Hom(p_a,U)$ and $\Hom(U,p_a)$. Their composition pairing takes values
in $\End(p_a)=K$ and is nondegenerate; invert its matrix to obtain dual
inclusions and projections. The sum of the resulting simple-block
projections is $\mathrm{id}_U$, by the complete semisimple decomposition
of \cite[Theorems~5.4.4--5.4.5]{Chen2014}. These identities hold over
$K$: they hold after extension to $\mathbb C$, and scalar extension is
faithful. Thus no search for an unspecified categorical decomposition
is required. Finite products of the factors and direct sums are handled
by tensor products and matrices of these maps.

It follows that every space
$\Hom(L,D^a_{ij}\otimes H)$ and every multiplication map defining $Z$
has computable structure constants over $K$. Let $Z_K$ denote this
finite algebra. Its trace weights belong to $K$, and matrix inversion
computes trace-dual bases, hence the coevaluation maps as well as the
evaluation maps. All terms and differentials of $B_i,D_i$ are now
explicit over $K$. Unnormalized planar bases suffice; the square-root
normalization used only in the proof of Lemma~\ref{sph:path-bound}
is not needed for this construction.

Given the input words $u,v$, form the word $w=u^{-1}vu$ and compute the
bounded projective complex $B_wP_{f,\mathbf1}$ by finite tensor products.
Cancel invertible scalar differential entries between equal shifted
projectives until none remain, as in Section~\ref{sec:layers}. Each
cancellation uses exact arithmetic in $K$ and removes two summands,
so the process terminates.

The positive-path-length ideal of $Z_K$ is nilpotent and its quotient is
a product of copies of $K$. It is therefore the radical, and it remains
the radical after extension to $\mathbb C$. The computed minimal complex
stays minimal under that extension. By uniqueness of minimal models,
it is equivalent over $\mathbb C$ to $P_{f,\mathbf1}$ precisely when its
only generator is that projective in placements $(0,0)$. This is a finite
check of its summands; when it succeeds there is no differential left.
Equation~\eqref{stab:single-detector} identifies success with $w\in A_X$,
equivalently $v\in uA_Xu^{-1}$. Taking $X=\varnothing$ decides whether a
word is the identity, and applying this test to the quotient of two words
decides their equality.
\end{proof}
